\documentclass[10pt,a4paper]{amsart}
\usepackage{amsmath,amssymb,amsthm,graphicx}
\usepackage{verbatim}
\newenvironment{reference}{\comment}{\endcomment}
%\newenvironment{reference}{}{}
\newenvironment{slogan}{\comment}{\endcomment}
\newenvironment{history}{\comment}{\endcomment}

% For commutative diagrams we use Xy-pic
\usepackage[all]{xy}

% We use 2cell for 2-commutative diagrams.
\xyoption{2cell}
% Theorem environments.
%
\theoremstyle{plain}
\newtheorem{theorem}[subsection]{Theorem}
\newtheorem{proposition}[subsection]{Proposition}
\newtheorem{lemma}[subsection]{Lemma}

\theoremstyle{definition}
\newtheorem{definition}[subsection]{Definition}
\newtheorem{example}[subsection]{Example}
\newtheorem{exercise}[subsection]{Exercise}
\newtheorem{situation}[subsection]{Situation}

\theoremstyle{remark}
\newtheorem{remark}[subsection]{Remark}
\newtheorem{remarks}[subsection]{Remarks}

\numberwithin{equation}{subsection}

% Macros
%
\def\lim{\mathop{\mathrm{lim}}\nolimits}
\def\colim{\mathop{\mathrm{colim}}\nolimits}
\def\Spec{\mathop{\mathrm{Spec}}}
\def\Hom{\mathop{\mathrm{Hom}}\nolimits}
\def\Ext{\mathop{\mathrm{Ext}}\nolimits}
\def\SheafHom{\mathop{\mathcal{H}\!\mathit{om}}\nolimits}
\def\SheafExt{\mathop{\mathcal{E}\!\mathit{xt}}\nolimits}
\def\Sch{\mathit{Sch}}
\def\Mor{\mathop{\mathrm{Mor}}\nolimits}
\def\Ob{\mathop{\mathrm{Ob}}\nolimits}
\def\Sh{\mathop{\mathit{Sh}}\nolimits}
\def\NL{\mathop{N\!L}\nolimits}
\def\CH{\mathop{\mathrm{CH}}\nolimits}
\def\proetale{{pro\text{-}\acute{e}tale}}
\def\etale{{\acute{e}tale}}
\def\QCoh{\mathit{QCoh}}
\def\Ker{\mathop{\mathrm{Ker}}}
\def\Im{\mathop{\mathrm{Im}}}
\def\Coker{\mathop{\mathrm{Coker}}}
\def\Coim{\mathop{\mathrm{Coim}}}

% Boxtimes
%
\DeclareMathSymbol{\boxtimes}{\mathbin}{AMSa}{"02}

%
% Macros for moduli stacks/spaces
%
\def\QCohstack{\mathcal{QC}\!\mathit{oh}}
\def\Cohstack{\mathcal{C}\!\mathit{oh}}
\def\Spacesstack{\mathcal{S}\!\mathit{paces}}
\def\Quotfunctor{\mathrm{Quot}}
\def\Hilbfunctor{\mathrm{Hilb}}
\def\Curvesstack{\mathcal{C}\!\mathit{urves}}
\def\Polarizedstack{\mathcal{P}\!\mathit{olarized}}
\def\Complexesstack{\mathcal{C}\!\mathit{omplexes}}
% \Pic is the operator that assigns to X its picard group, usage \Pic(X)
% \Picardstack_{X/B} denotes the Picard stack of X over B
% \Picardfunctor_{X/B} denotes the Picard functor of X over B
\def\Pic{\mathop{\mathrm{Pic}}\nolimits}
\def\Picardstack{\mathcal{P}\!\mathit{ic}}
\def\Picardfunctor{\mathrm{Pic}}
\def\Deformationcategory{\mathcal{D}\!\mathit{ef}}
\title{A unique prime over p plus going up (or down) pins down the localisation at p}
\author{The Stacks Project Authors}
\date{stacks-project:algebra.tex:local-snapshot-20261012; GNU FDL 1.2 or later}
\begin{document}
\maketitle
\noindent\textit{Editorial note.} All delivered mathematical text is the author's own, copied from the pinned Stacks Project chapter file algebra.tex; no statement or proof step was rewritten, reordered or removed.\par\medskip
\noindent\textit{Editorial note.} Preamble provenance, as required for this source: the frozen snapshot stages the chapter's own preamble file beside it (preamble.tex, 217 lines, 8116 bytes, the same shared Stacks-wide preamble the published chow and etale-cohomology cards use); the first line of the pinned chapter file is \textbackslash{}input\{preamble\}, so no preamble text lives inside the chapter itself. The card takes that file's lines 8-18 (the verbatim package and the three \textbackslash{}comment-based environments reference/slogan/history at lines 9-12, section 13.2, together with the Xy-pic setup) and its lines 154 to the runtime end of the file (the native theorem declarations and mathematical macros). Not one character was changed, no package was added, no macro was defined, and the excerpt serves only as preamble, never as a span coordinate.\par\medskip
\noindent\textit{Editorial note.} The same-chapter pointers are located at their \textbackslash{}label lines inside the pinned chapter file; the cross-chapter pointers are located in the project's own tag list staged in the snapshot (sources/stacks-chow/tags.txt), at the verbatim lines carrying their canonical tags.\par\medskip
\noindent\textit{Editorial note.} The counted spans contain no \textbackslash{}ref and no \textbackslash{}cite, so the card carries no navigation stub and no bibliography.\par\medskip
\noindent\textit{Allowed standard prerequisites (editorial).}\begin{itemize}
\item Going up and going down as lifting properties for chains of primes. Used in the two branches of the counted proof (the diagram chase in lines 10035-10058). Carries a chain of primes in the base to a chain in the source, which is what forces the primes above p to be unique.
\item Primes of a localisation correspond to primes below the localised prime. Used at the start of the counted proof (lines 10031-10034). Describes the points of the localisation as the primes whose contraction lies below p.
\end{itemize}\medskip
\begin{lemma}
\label{lemma-unique-prime-over-localize-below}
Let $R \to S$ be a ring map.
Let $\mathfrak p \subset R$ be a prime.
Assume that
\begin{enumerate}
\item there exists a unique prime $\mathfrak q \subset S$ lying over
$\mathfrak p$, and
\item either
\begin{enumerate}
\item going up holds for $R \to S$, or
\item going down holds for $R \to S$ and there is at most one prime
of $S$ above every prime of $R$.
\end{enumerate}
\end{enumerate}
Then $S_{\mathfrak p} = S_{\mathfrak q}$.
\end{lemma}

\begin{proof}
Consider any prime $\mathfrak q' \subset S$ which corresponds to
a point of $\Spec(S_{\mathfrak p})$. This means that
$\mathfrak p' = R \cap \mathfrak q'$ is contained in $\mathfrak p$.
Here is a picture
$$
\xymatrix{
\mathfrak q' \ar@{-}[d] \ar@{-}[r] & ? \ar@{-}[r] \ar@{-}[d] & S \ar@{-}[d] \\
\mathfrak p' \ar@{-}[r] & \mathfrak p \ar@{-}[r] & R
}
$$
Assume (1) and (2)(a).
By going up there exists a prime $\mathfrak q'' \subset S$
with $\mathfrak q' \subset \mathfrak q''$ and $\mathfrak q''$
lying over $\mathfrak p$. By the uniqueness of $\mathfrak q$ we
conclude that $\mathfrak q'' = \mathfrak q$. In other words
$\mathfrak q'$ defines a point of $\Spec(S_{\mathfrak q})$.

\medskip\noindent
Assume (1) and (2)(b).
By going down there exists a prime $\mathfrak q'' \subset \mathfrak q$
lying over $\mathfrak p'$. By the uniqueness of primes lying over
$\mathfrak p'$ we see that $\mathfrak q' = \mathfrak q''$.  In other words
$\mathfrak q'$ defines a point of $\Spec(S_{\mathfrak q})$.

\medskip\noindent
In both cases we conclude that the map
$\Spec(S_{\mathfrak q}) \to \Spec(S_{\mathfrak p})$
is bijective. Clearly this means all the elements of $S - \mathfrak q$
are all invertible in $S_{\mathfrak p}$, in other words
$S_{\mathfrak p} = S_{\mathfrak q}$.
\end{proof}

\end{document}
